\section{Finite game and realization-plan formulation}
Write a complete own path as $i=(i_0,\ldots,i_{T-1})$ and the opponent path
as $j=(j_0,\ldots,j_{T-1})$. Both use a state-independent menu of size $b$.
At an own update $t$, the information available for a new choice is exactly
$(i_{<t},j_{<t})$. All previously selected own commands have then executed;
their contents are already in $i_{<t}$. The choice is an element of
$\{0,\ldots,b-1\}^{e-t}$, where $e$ is the next own update or the horizon.
The opponent's currently outstanding suffix is excluded. These information
sets have perfect recall: every earlier own choice and every earlier
observation is recoverable from the present prefix.

For each own information set $I$, let $\sigma(I)$ be its incoming own
sequence and let $\sigma(I)a$ denote each block-extended sequence. Include
an empty sequence with realization weight one. A realization plan satisfies
\[
 x_{\emptyset}=1,\qquad x\ge0,\qquad
 \sum_{a\in A(I)}x_{\sigma(I)a}=x_{\sigma(I)}.
\]
Collect these equations as $Ex=e$; write the opponent equations as $Fy=f$.
For every terminal pair $(i,j)$ put its full trajectory payoff into the
matrix entry indexed by its final own and opponent sequences. If more than
one terminal node shares an index, add its chance-weighted contribution.
There are no exogenous chance moves in the formation benchmark. The expected
payoff is $x^\top Py$ and the value is
\[
 \max_{x\ge0,Ex=e}\min_{y\ge0,Fy=f}x^\top Py.
\]
The inner problem has dual $\max_v f^\top v$ subject to $F^\top v\le P^\top x$.
Consequently the maximizer solves
\begin{equation}
 \max_{x,v}f^\top v\quad\text{s.t.}\quad
 Ex=e,\ x\ge0,\ F^\top v\le P^\top x,
 \label{eq:sf}
\end{equation}
with unrestricted $v$. Finite game duality and perfect-recall realization
equivalence justify this standard formulation \citep{stengel1996}.
Its size is linear in the represented game tree, which can itself grow
exponentially in $T$.

For a numerical candidate $(x,y)$, nonnegative conditional child weights
are normalized at every information set and multiplied by repaired parent
reach. If proposed child mass is zero, use a uniform conditional distribution.
This constructs feasible realization plans in exact arithmetic. Report the
actual flow residual and the maximum repair magnitude. Two unrestricted
responses then give
\[
 \min_{Fy'=f,y'\ge0}x^\top Py'\ \le V\ \le
 \max_{Ex'=e,x'\ge0}{x'}^\top Py.
\]
Feasibility repair does not correct a wrongly specified information set.
The independent tiny-game normal-form enumeration is therefore a separate
test of model representation.

\section{Exact counterfactual pricing}
Fix a causal opponent policy $q$. For each externally specified own path
$i$, let $C_q(j\mid i)$ be the probability of opponent path $j$. Causality
means that for any $t$ and opponent prefix $j_{<t}$,
\begin{equation}
 \sum_{j_{\ge t}}C_q(j\mid i)
\end{equation}
depends on $i_{<t}$ and not on any later own action. This follows by multiplying
the opponent's conditional decision probabilities, or by averaging causal
pure policies. An opponent may retain a private command suffix inside a
block; the marginalization integrates over that suffix rather than revealing it.

Initialize $H_T(i,j)=A(i,j)C_q(j\mid i)$. If $t<e$ are successive own
boundaries, recursively define
\begin{equation}
 H_t(i_{<t},j_{<t})=max_{a\in\{0,\ldots,b-1\}^{e-t}}
 \sum_{z\in\{0,\ldots,b-1\}^{e-t}}
 H_e(i_{<t}a,j_{<t}z).
 \label{eq:price}
\end{equation}
The terminal entries are weighted by opponent realization probabilities,
not by an own policy. Suppose all decisions from $e$ onward have been optimized.
At time $t$, the own block $a$ must be fixed before $z$ is observed. Thus $z$
is summed for each candidate $a$; the maximizing block is then selected.
The probability of the already executed opponent prefix is independent of
this candidate. It can remain as an unnormalized common factor throughout
the recurrence. The induction hypothesis permits different optimal
continuations at distinct newly observed prefixes at $e$. This proves
\eqref{eq:price} and the root best-response value. Any stored maximizing
choice at a zero-weight prefix is admissible.

For $b\ge2$, the number of weighted entries processed is at most
\[
 b^{2T}+b^{2(T-1)}+\cdots+1\le\frac{b^{2T}}{1-b^{-2}}.
\]
The payoff matrix and response arrays therefore take $O(b^{2T})$ arithmetic
and memory. Constructing the physical payoff and repeating this procedure
inside an equilibrium solver remain additional costs. This is a specialized
history-indexed response calculation, not a polynomial algorithm in the
compact physical-state description.

\section{Restricted policy generation and stopping}
Against a committed opponent, a feasible pure own policy $p_m$ induces one
own path $i_m(j)$ for each complete opponent path $j$. Form a restricted
matrix with entries $B_{mj}=A(i_m(j),j)$, retaining every opponent column.
Let $w$ and $q$ be restricted equilibrium distributions. The mixed own policy
has full-game lower guarantee
\[
 L=\min_j\sum_m w_mB_{mj},
\]
and exact pricing against $q$ gives $U=\BR_S(q)\ge V(S,C)$. Hence $[L,U]$
is valid. If $U=L$ in exact arithmetic, the restricted mixture and $q$ are
an equilibrium of the complete game. If $U>L$, a maximizing response cannot
be a previously available row, since all old rows have payoff at most the
restricted value against $q$. Adding that response strictly enlarges the
finite set of available policies. The exact procedure terminates after
finitely many additions, but this worst-case count can be exponential.
This is the standard response-oracle argument
\citep{mcmahan2003,bosansky2014}.

The implementation attempts at most 64 policy additions before switching to
the complete realization-plan LP for the same game. The first attempt's
time is included. A duplicate response with a numerical gap is reported as
a stall, not used as a proof of convergence. All positive mixture weights
enter full-response evaluation; pruning tiny positive weights would require
an additional error budget.

\section{Coarsening a full-update behavioral strategy}
Let $\pi_u(a\mid i_{<u},j_{<u})$ be a full-update source strategy, completed
by uniform distributions at its zero-reach histories. Fix one legal reference
opponent action $r_\star$, the straight action in the experiment. For a target
block $[t,e)$ and actual observed prefixes $(i_{<t},j_{<t})$, define
\begin{equation}
 \widetilde\pi_{t:e}(a_t,\ldots,a_{e-1}\mid i_{<t},j_{<t})
 =\prod_{u=t}^{e-1}
 \pi_u\left(a_u\mid i_{<t}a_{t:u},j_{<t}r_\star^{u-t}\right).
 \label{eq:coarse}
\end{equation}
For each fixed earlier own block prefix, the last factor sums to one over
$a_{e-1}$. Repeating this argument backward proves normalization of
\eqref{eq:coarse}. Its factors use only observed prefixes, private sampled
own commands and a deterministic hypothetical continuation. No actual future
opponent command is supplied. State-independent menus make all resulting
blocks legal. Thus the construction is a feasible target-calendar policy.

The source need not be optimal for this feasibility result. If it is obtained
as a root mixture, first convert its realization distribution to behavioral
kernels using perfect recall. Deleting updates after this conversion defines
a new behavioral strategy; it need not preserve the source's mixture over
latent pure plans or its outcome distribution. This distinction matters at
histories that become newly reachable after coarsening.

Let $U_F$ be any valid upper bound on the full-update value for the same
opponent class. A full response to $\widetilde\pi$ gives $L'$. Policy inclusion
and feasibility imply
\[
 0\le V(F,R)-V(S',R)\le U_F-L'.
\]
To evaluate the minimizing response, apply \eqref{eq:price} to the inverted
matrix $-A^\top$ weighted by the counterfactual realization law of
$\widetilde\pi$, using the opponent calendar $R$. This directly quantifies
the actual projected policy. It avoids unproved smoothness assumptions but
can yield a loose bound. It does not strengthen the general abstraction
decomposition of \citet{kroer2018}.

\section{Calendar search and transferable opponents}
A search node consists of included epochs $I$ and undecided epochs $D$.
Excluded epochs are the complement of $I\cup D$. Its remaining family is
\[
 \mathcal C(I,D)=\{S:I\subseteq S\subseteq I\cup D,\ |S|\le K\}.
\]
For any $S$ in this family, $V(S,R)\le V(I\cup D,R)$ by policy inclusion.
The latter game is a valid relaxation even if $|I\cup D|>K$. Its numerical
upper bound is therefore a node upper bound. A fixed opponent strategy
remains feasible at every node because its own calendar, available history
and action menu are unchanged. Thus $\BR_{I\cup D}(q)$ is another valid upper
bound, and a minimum over such bounds is safe.

If an upper bound is at most the incumbent feasible lower bound plus
$\varepsilon$, prune the node while retaining that bound in the terminal
bound list. Otherwise split on one undecided epoch, producing the included
and excluded child. These children partition the remaining family. When
$|I|=K$, only $I$ remains. When $|I\cup D|\le K$, inclusion implies that
$I\cup D$ dominates all calendars within the node, so solve only that calendar.
After finite exhaustion, every feasible calendar has either been evaluated
or belongs to a represented pruned node. The largest retained upper bound,
together with the best feasible lower bound, encloses $W_K$. If unresolved
inner solves leave a gap above $\varepsilon$, the search result is unresolved.
Finite termination alone does not establish numerical optimality.

\section{Interval losses and the shortest-path construction}
For opponent $F$, every physical epoch is a potential common boundary. The
fully replanning continuation tree supplies values $v_t(h)$ at every public
history. Let $b_{t,e}(h)$ denote the value of a block game with Blue committed,
Red fully replanning and terminal continuation $v_e$. Define
$\delta_{t,e}(h)=v_t(h)-b_{t,e}(h)$ and
$d(t,e)=\max_h\delta_{t,e}(h)$. These deficits are nonnegative because the
block restriction only removes own decisions. Adjacent boundaries give the
identical one-stage simultaneous game, hence $d(t,t+1)=0$.

Fix a calendar and stitch the maximizing block policies for its successive
intervals. Write $G_t(h)$ for the lower guarantee of this stitched policy
from $t$. If $G_e(h')\ge v_e(h')-D_e$ for every possible successor history,
the same block policy guarantees
\[
 G_t(h)\ge b_{t,e}(h)-D_e
       =v_t(h)-\delta_{t,e}(h)-D_e
       \ge v_t(h)-d(t,e)-D_e.
\]
The first inequality holds against every within-block Red policy, including
one chosen with full knowledge of the public calendar. Future Red choices
are covered by the uniform continuation guarantee. Because Red updates at
every epoch and Blue's block ends at $e$, no private command crosses this
boundary. Starting with $D_T=0$ and inducting backward proves the main
interval-loss theorem. This also constructs a feasible policy attaining
the bound in exact arithmetic.

For nonnegative certified edge bounds $\bar d(t,e)\ge d(t,e)$, set
$D(0,0)=0$ and all other $D(0,e)=+\infty$. The recurrence
\[
 D(j,e)=\min_{0\le t<e}\{D(j-1,t)+\bar d(t,e)\}
\]
uses exactly $j$ intervals. Its budget answer is
$\min_{1\le j\le K+1}D(j,T)$. The minimum over all smaller $j$ is essential:
the edge upper bounds need not satisfy a triangle inequality. Storing each
minimizing predecessor recovers a calendar. The resulting value lower bound
is $L_F-\min_{j\le K+1}D(j,T)$ when $L_F\le V(F,F)$. It need not be the
optimal calendar value, and a sufficient retention budget from this bound
need not be the true minimum budget.

For numerical edge construction, suppose the full tree uses approximations
$\widehat v_t$ with uniform errors $\epsilon_t$. A lower bound
$\widehat L_{t,e}(h)$ for the block game with terminal reward
$\widehat v_e$ gives
\[
 d(t,e)\le\max_h\{
 \widehat v_t(h)+\epsilon_t-\widehat L_{t,e}(h)+\epsilon_e\}.
\]
The implemented upper bound also takes its maximum with zero. Adjacent
edges use their exact identity $d=0$. Every local game and the full tree
are charged as preprocessing; the cheap graph search must not be reported
as the total method cost. The finite-history maximum is deliberately
conservative and requires no global smoothness of the directional kernel.

\section{A sufficient modular special case}
Suppose there are independent component games indexed by potential update
epochs $m$. Component $m$ has maximizing pure-policy set $P_m^0$ without its
update and $P_m^1\supseteq P_m^0$ with its update, a fixed minimizing set $Q_m$,
and payoff $a_m(p_m,q_m)$. Assume the whole game's pure-policy sets are exactly
the Cartesian products of the component sets and its payoff is the sum of
the component payoffs. These are structural assumptions: there are no shared
opponent resources, inter-component physical state constraints or additional
cross-component observations that enlarge these sets.

Let $v_m^d$ be the minimax value of component $m$ under $d\in\{0,1\}$. The
product of component maximizing equilibrium policies guarantees
$\sum_m v_m^{\mathbf1\{m\in S\}}$ against every joint minimizing policy,
because each of its component marginals is feasible. The product of component
minimizing equilibria gives the matching upper bound against every joint
maximizing policy. Therefore
\[
 V(S)=\sum_m v_m^0+\sum_{m\in S}(v_m^1-v_m^0),
\]
which is modular. The assumptions define a checkable special class and are
not claimed for the formation benchmark. The shared trajectory and target
history prevent this product factorization in general.

\section{Boundary conditions and historical counterexamples}
The angular kernel changes firing sector at specified bearings. A nonzero
one-sided jump is incompatible with a finite global Lipschitz constant:
the ratio of the jump to an arbitrarily small bearing separation is unbounded.
All current certificates operate on the explicitly evaluated finite payoff
matrix, including its boundary convention. They do not certify angular-grid
refinement, continuous-control limits or the physical accuracy of the kernel.
The archived floating-boundary amendment groups values within $10^{-8}$
degrees at the designated $\pm30$ and $\pm150$ degree boundaries. No payment
matrix is projected onto an antisymmetric matrix to pass a mirror test.

Exchange oddness is not a Bellman collapse. Take absorbing states $s=\pm1$
with no actions and stage reward $s$. Exchanging players sends $s$ to $-s$,
but the two-step value is $2s$, not the one-stage reward. A separate issue is
pure versus mixed play: the rock--paper--scissors antisymmetric matrix has
pure max--min value $-1$ and mixed minimax value zero. Zero-value claims must
identify the correct exchange-fixed state and the matching policy classes.

Changing a forced operating speed changes the feasible trajectories and is
not a nested enlargement of the action menu. Even an actually nested ability
set guarantees only that each value with fixed opposing class cannot decrease;
it does not make the difference of two such values monotone. Archived v12
range and v13 mobility contrasts therefore remain separate negative evidence,
not premises of the present budget theorem.

\section{Numerical intervals and decision logic}
The scientific objective is the finite deterministic game. Declared numerical
intervals come from feasible strategies, complete responses, local error
propagation and explicit floating-point padding. They are not interval-arithmetic
proofs: arithmetic and physical-payoff rounding are not enclosed with directed
rounding. The implementation reports this scope with every solver result.

For adaptation retention $\rho$, define
$M_K=W_K-\rho V_F-(1-\rho)V_0$. Given value intervals, the lower and upper
margins are
\[
 M_K^- = L_K-\rho U_F-(1-\rho)U_0,\qquad
 M_K^+ = U_K-\rho L_F-(1-\rho)L_0.
\]
A nonnegative lower margin proves sufficiency, and a negative upper margin
proves insufficiency. The smallest possibly sufficient budget and the smallest
certified sufficient budget are reported separately. Equality of these budgets
certifies minimality. When the adaptation interval is too close to zero, the
reported quantity is absolute loss; no unstable retention percentage is formed.
These are optimization intervals on a deterministic finite design, not sampling
confidence intervals.
